\begin{center}\bfseries Abstract\end{center}
\begin{narrowtext}
Sub-Neptune exoplanets represent the most abundant class of planets detected in the known
exoplanet population. Despite their ubiquity, their fundamental properties remain poorly
characterized. In particular, first-principles simulations suggest that hydrogen and silicate
are completely miscible at the magma-atmosphere boundary of many young sub-Neptunes, so that
hydrogen exsolves as the planet cools, but the timing of this exsolution remains largely
unexplored.\par
\vspace{3pt}
In this work, we investigate the thermodynamic limits on hydrogen exsolution from
H$_2$--MgSiO$_3$ mixtures using \Atm{}, a Python-based framework designed for coupled
atmosphere-interior modeling. We implement the final 2026 mixing model with a
composition-consistent basis, compare the equilibrium onset of exsolution (the binodal) with
the loss of local stability (the spinodal), and recalibrate four smooth model classes against
the source simulations.\par
\vspace{3pt}
With the published fit, a parcel with 3\,wt\% hydrogen at 2\,GPa enters the miscibility gap at
\pubTbTwo{}\,K but becomes unstable only at \pubTspTwo{}\,K; if nucleation is suppressed, up
to \pubFhTwo{}\% of its hydrogen is withheld over this \pubIntTwo{}\,K interval. This delay is
not constrained by the source data: refits of the same data give median intervals of
\IntRangeLo{} to \IntRangeHi{}\,K. Equilibrium miscibility data therefore do not determine when
hydrogen separates. A direct measurement of the composition curvature would fix the spinodal
to within about \curvSdPooled{}\,K.
\end{narrowtext}
\vspace{10pt}
\noindent\rule{\textwidth}{0.6pt}
\vspace{10pt}
\begin{center}\bfseries Samenvatting\end{center}
\begin{narrowtext}
Sub-Neptunus-exoplaneten vormen de meest talrijke klasse binnen de huidige
exoplaneetpopulatie. Ondanks hun alomtegenwoordigheid blijven hun fundamentele
eigenschappen echter slecht gekarakteriseerd. Simulaties vanuit eerste principes suggereren
dat waterstof en silicaat volledig mengbaar zijn aan de grens tussen magma en atmosfeer van
veel jonge sub-Neptunussen. Waterstof ontmengt dan wanneer de planeet afkoelt, maar het
tijdstip van deze ontmenging is tot op heden nauwelijks onderzocht.\par
\vspace{3pt}
In dit werk bestuderen wij de thermodynamische grenzen van de ontmenging van waterstof uit
H$_2$--MgSiO$_3$-mengsels met behulp van \Atm{}, een in Python ontwikkeld raamwerk voor
gekoppelde atmosfeer-mantelmodellering. Wij implementeren het definitieve mengmodel uit 2026
met een consistente samenstellingsbasis, vergelijken het begin van de ontmenging in evenwicht
(de binodaal) met het verlies van lokale stabiliteit (de spinodaal), en kalibreren vier gladde
modelklassen opnieuw op de oorspronkelijke simulaties.\par
\vspace{3pt}
Met de gepubliceerde fit komt een mengsel met 3\,massa\% waterstof bij 2\,GPa de
mengbaarheidskloof binnen bij \pubTbTwo{}\,K, maar wordt het pas instabiel bij
\pubTspTwo{}\,K; als kiemvorming onderdrukt wordt, blijft tot \pubFhTwo{}\% van de waterstof
achter over dit interval van \pubIntTwo{}\,K. De oorspronkelijke gegevens leggen deze
vertraging niet vast: nieuwe fits aan dezelfde gegevens geven mediane intervallen van
\IntRangeLo{} tot \IntRangeHi{}\,K. Evenwichtsgegevens over mengbaarheid bepalen dus niet
wanneer waterstof ontmengt.
\end{narrowtext}
